
\documentclass{article}
\usepackage{colortbl}
\usepackage{makecell}
\usepackage{multirow}
\usepackage{supertabular}

\begin{document}

\newcounter{utterance}

\centering \large Interaction Transcript for game `hot\_air\_balloon', experiment `air\_balloon\_survival\_de\_complexity\_easy', episode 3 with deepseek{-}chat{-}v3.1{-}t1.0.
\vspace{24pt}

{ \footnotesize  \setcounter{utterance}{1}
\setlength{\tabcolsep}{0pt}
\begin{supertabular}{c@{$\;$}|p{.15\linewidth}@{}p{.15\linewidth}p{.15\linewidth}p{.15\linewidth}p{.15\linewidth}p{.15\linewidth}}
    \# & $\;$A & \multicolumn{4}{c}{Game Master} & $\;\:$B\\
    \hline

    \theutterance \stepcounter{utterance}  
    & & \multicolumn{4}{p{0.6\linewidth}}{
        \cellcolor[rgb]{0.9,0.9,0.9}{
            \makecell[{{p{\linewidth}}}]{
                \texttt{\tiny{[P1$\langle$GM]}}
                \texttt{Du nimmst an einem kollaborativen Verhandlungspiel Teil.} \\
\\ 
\texttt{Zusammen mit einem anderen Teilnehmer musst du dich auf eine Reihe von Gegenständen entscheiden, die behalten werden. Jeder von euch hat eine persönliche Verteilung über die Wichtigkeit der einzelnen Gegenstände. Jeder von euch hat eine eigene Meinung darüber, wie wichtig jeder einzelne Gegenstand ist (Gegenstandswichtigkeit). Du kennst die Wichtigkeitsverteilung des anderen Spielers nicht. Zusätzlich siehst du, wie viel Aufwand jeder Gegenstand verursacht.  } \\
\texttt{Ihr dürft euch nur auf eine Reihe von Gegenständen einigen, wenn der Gesamtaufwand der ausgewählten Gegenstände den Maximalaufwand nicht überschreitet:} \\
\\ 
\texttt{Maximalaufwand = 8770} \\
\\ 
\texttt{Hier sind die einzelnen Aufwände der Gegenstände:} \\
\\ 
\texttt{Aufwand der Gegenstände = \{"C32": 830, "C52": 638, "B54": 607, "C16": 106, "C50": 630, "A08": 382, "C59": 773, "B51": 455, "B55": 262, "B59": 634, "A77": 60, "C14": 54, "C24": 849, "B24": 801, "C31": 324, "A52": 164, "C61": 832, "B79": 136, "C54": 980, "B56": 946, "A40": 645, "B11": 842, "B60": 106, "A81": 116, "A95": 875, "B38": 446, "C51": 649, "C03": 601, "A91": 252, "A93": 763, "C28": 213, "C05": 517, "B75": 520, "B96": 407, "B86": 125\}} \\
\\ 
\texttt{Hier ist deine persönliche Verteilung der Wichtigkeit der einzelnen Gegenstände:} \\
\\ 
\texttt{Werte der Gegenstandswichtigkeit = \{"C32": 138, "C52": 583, "B54": 868, "C16": 822, "C50": 783, "A08": 65, "C59": 262, "B51": 121, "B55": 508, "B59": 780, "A77": 461, "C14": 484, "C24": 668, "B24": 389, "C31": 808, "A52": 215, "C61": 97, "B79": 500, "C54": 30, "B56": 915, "A40": 856, "B11": 400, "B60": 444, "A81": 623, "A95": 781, "B38": 786, "C51": 3, "C03": 713, "A91": 457, "A93": 273, "C28": 739, "C05": 822, "B75": 235, "B96": 606, "B86": 968\}} \\
\\ 
\texttt{Ziel:} \\
\\ 
\texttt{Dein Ziel ist es, eine Reihe von Gegenständen auszuhandeln, die dir möglichst viel bringt (d. h. Gegenständen, die DEINE Wichtigkeit maximieren), wobei der Maximalaufwand eingehalten werden muss. Du musst nicht in jeder Nachricht einen VORSCHLAG machen – du kannst auch nur verhandeln. Alle Taktiken sind erlaubt!} \\
\\ 
\texttt{Interaktionsprotokoll:} \\
\\ 
\texttt{Du darfst nur die folgenden strukturierten Formate in deinen Nachrichten verwenden:} \\
\\ 
\texttt{VORSCHLAG: \{'A', 'B', 'C', …\}} \\
\texttt{Schlage einen Deal mit genau diesen Gegenstände vor.} \\
\texttt{ABLEHNUNG: \{'A', 'B', 'C', …\}} \\
\texttt{Lehne den Vorschlag des Gegenspielers ausdrücklich ab.} \\
\texttt{ARGUMENT: \{'...'\}} \\
\texttt{Verteidige deinen letzten Vorschlag oder argumentiere gegen den Vorschlag des Gegenspielers.} \\
\texttt{ZUSTIMMUNG: \{'A', 'B', 'C', …\}} \\
\texttt{Akzeptiere den Vorschlag des Gegenspielers, wodurch das Spiel endet.} \\
\texttt{STRATEGISCHE ÜBERLEGUNGEN: \{'...'\}} \\
\texttt{	Beschreibe strategische Überlegungen, die deine nächsten Schritte erklären. Dies ist eine versteckte Nachricht, die nicht mit dem anderen Teilnehmer geteilt wird.} \\
\\ 
\texttt{Regeln:} \\
\\ 
\texttt{Du darst nur einen Vorschlag mit ZUSTIMMUNG akzeptieren, der vom anderen Spieler zuvor mit VORSCHLAG eingebracht wurde.} \\
\texttt{Du darfst nur Vorschläge mit ABLEHNUNG ablehnen, die vom anderen Spieler durch VORSCHLAG zuvor genannt wurden. } \\
\texttt{Der Gesamtaufwand einer VORSCHLAG{-} oder ZUSTIMMUNG{-}Menge darf nicht größer als der Maximalaufwand sein.  } \\
\texttt{Offenbare deine versteckte Wichtigkeitsverteilung nicht.} \\
\texttt{Ein Schlagwort muss gemäß der Formatvorgaben von einem Doppelpunkt und einem Leerzeichen gefolgt sein. Das Argument muss eine Python{-}Menge mit 0 oder mehr Strings sein.  } \\
\texttt{Die Form muss also SCHLAGWORT: \{...\} sein.} \\
\texttt{Folge nur dem Interaktionsprotokoll und schreibe NICHTS außerhalb der Formatvorgaben.} \\
\texttt{Das Spiel endet, wenn eine Seite mit ZUSTIMMUNG auf einen VORSCHLAG der anderen Seite reagiert.  } \\
\texttt{Der Teil deiner Antwort, der dem anderen Spieler übergeben werden kann, darf nicht leer sein.  } \\
\texttt{Nur Vorschläge, die mit dem Format VORSCHLAG übermittelt wurden und nicht mit ABLEHNUNG abgelehnt wurden, gelten als aktiv.  } \\
\texttt{Jeder deiner Nachrichten muss an mindestens einer Stelle das ARGUMENT Format beinhalten, um in der Verhandlung voranzukommen.} \\
\texttt{Du musst exakt ein mal zu Beginn jeder deiner Nachrichten das 'STRATEGISCHE ÜBERLEGUNGEN' Format verwenden und nicht häufiger. Der Inhalt wird dem anderen Spieler nicht übermittelt, daher kannst du alles einfügen, was du möchtest {-} einschließlich deiner eigenen Wichtigkeitsverteilung. Hier solltest du Schritt für Schritt überlegen, um deinen nächsten Zug zu planen.} \\
            }
        }
    }
    & \\ \\

    \theutterance \stepcounter{utterance}  
    & \multicolumn{4}{p{0.6\linewidth}}{
        \cellcolor[rgb]{0.9,0.9,0.9}{
            \makecell[{{p{\linewidth}}}]{
                \texttt{\tiny{[P1$\rangle$GM]}}
                \texttt{STRATEGISCHE ÜBERLEGUNGEN: \{'Ich muss einen Vorschlag machen, der meinen Gesamtnutzen maximiert, während der Maximalaufwand von 8770 eingehalten wird. Meine Top{-}Gegenstände mit hohem Wert und niedrigem Aufwand sind: B86 (Wert 968, Aufwand 125), C16 (822/106), C05 (822/517), B54 (868/607), A40 (856/645), B56 (915/946), B59 (780/634), C31 (808/324). Ich sollte versuchen, viele dieser wertvollen Gegenstände in meinen Vorschlag aufzunehmen, während ich weniger wertvolle Gegenstände opfere, um Platz für Verhandlungen zu lassen.'\}} \\
\\ 
\texttt{VORSCHLAG: \{'B86', 'C16', 'C05', 'B54', 'A40', 'B59', 'C31', 'A81', 'B60', 'C14', 'A77', 'B79', 'C28', 'A52', 'B55'\}} \\
\\ 
\texttt{ARGUMENT: \{'Dieser Vorschlag beinhaltet mehrere Gegenstände mit gutem Wert{-}Aufwand{-}Verhältnis und lässt Raum für Anpassungen. Ich bin offen für Diskussionen über bestimmte Gegenstände, um eine für beide Seiten vorteilhafte Lösung zu finden.'\}} \\
            }
        }
    }
    & & \\ \\

    \theutterance \stepcounter{utterance}  
    & & \multicolumn{4}{p{0.6\linewidth}}{
        \cellcolor[rgb]{0.9,0.9,0.9}{
            \makecell[{{p{\linewidth}}}]{
                \texttt{\tiny{[GM$\rangle$P2]}}
                \texttt{Du nimmst an einem kollaborativen Verhandlungspiel Teil.} \\
\\ 
\texttt{Zusammen mit einem anderen Teilnehmer musst du dich auf eine Reihe von Gegenständen entscheiden, die behalten werden. Jeder von euch hat eine persönliche Verteilung über die Wichtigkeit der einzelnen Gegenstände. Jeder von euch hat eine eigene Meinung darüber, wie wichtig jeder einzelne Gegenstand ist (Gegenstandswichtigkeit). Du kennst die Wichtigkeitsverteilung des anderen Spielers nicht. Zusätzlich siehst du, wie viel Aufwand jeder Gegenstand verursacht.  } \\
\texttt{Ihr dürft euch nur auf eine Reihe von Gegenständen einigen, wenn der Gesamtaufwand der ausgewählten Gegenstände den Maximalaufwand nicht überschreitet:} \\
\\ 
\texttt{Maximalaufwand = 8770} \\
\\ 
\texttt{Hier sind die einzelnen Aufwände der Gegenstände:} \\
\\ 
\texttt{Aufwand der Gegenstände = \{"C32": 830, "C52": 638, "B54": 607, "C16": 106, "C50": 630, "A08": 382, "C59": 773, "B51": 455, "B55": 262, "B59": 634, "A77": 60, "C14": 54, "C24": 849, "B24": 801, "C31": 324, "A52": 164, "C61": 832, "B79": 136, "C54": 980, "B56": 946, "A40": 645, "B11": 842, "B60": 106, "A81": 116, "A95": 875, "B38": 446, "C51": 649, "C03": 601, "A91": 252, "A93": 763, "C28": 213, "C05": 517, "B75": 520, "B96": 407, "B86": 125\}} \\
\\ 
\texttt{Hier ist deine persönliche Verteilung der Wichtigkeit der einzelnen Gegenstände:} \\
\\ 
\texttt{Werte der Gegenstandswichtigkeit = \{"C32": 138, "C52": 583, "B54": 868, "C16": 822, "C50": 783, "A08": 65, "C59": 262, "B51": 121, "B55": 508, "B59": 780, "A77": 461, "C14": 484, "C24": 668, "B24": 389, "C31": 808, "A52": 215, "C61": 97, "B79": 500, "C54": 30, "B56": 915, "A40": 856, "B11": 400, "B60": 444, "A81": 623, "A95": 781, "B38": 786, "C51": 3, "C03": 713, "A91": 457, "A93": 273, "C28": 739, "C05": 822, "B75": 235, "B96": 606, "B86": 968\}} \\
\\ 
\texttt{Ziel:} \\
\\ 
\texttt{Dein Ziel ist es, eine Reihe von Gegenständen auszuhandeln, die dir möglichst viel bringt (d. h. Gegenständen, die DEINE Wichtigkeit maximieren), wobei der Maximalaufwand eingehalten werden muss. Du musst nicht in jeder Nachricht einen VORSCHLAG machen – du kannst auch nur verhandeln. Alle Taktiken sind erlaubt!} \\
\\ 
\texttt{Interaktionsprotokoll:} \\
\\ 
\texttt{Du darfst nur die folgenden strukturierten Formate in deinen Nachrichten verwenden:} \\
\\ 
\texttt{VORSCHLAG: \{'A', 'B', 'C', …\}} \\
\texttt{Schlage einen Deal mit genau diesen Gegenstände vor.} \\
\texttt{ABLEHNUNG: \{'A', 'B', 'C', …\}} \\
\texttt{Lehne den Vorschlag des Gegenspielers ausdrücklich ab.} \\
\texttt{ARGUMENT: \{'...'\}} \\
\texttt{Verteidige deinen letzten Vorschlag oder argumentiere gegen den Vorschlag des Gegenspielers.} \\
\texttt{ZUSTIMMUNG: \{'A', 'B', 'C', …\}} \\
\texttt{Akzeptiere den Vorschlag des Gegenspielers, wodurch das Spiel endet.} \\
\texttt{STRATEGISCHE ÜBERLEGUNGEN: \{'...'\}} \\
\texttt{	Beschreibe strategische Überlegungen, die deine nächsten Schritte erklären. Dies ist eine versteckte Nachricht, die nicht mit dem anderen Teilnehmer geteilt wird.} \\
\\ 
\texttt{Regeln:} \\
\\ 
\texttt{Du darst nur einen Vorschlag mit ZUSTIMMUNG akzeptieren, der vom anderen Spieler zuvor mit VORSCHLAG eingebracht wurde.} \\
\texttt{Du darfst nur Vorschläge mit ABLEHNUNG ablehnen, die vom anderen Spieler durch VORSCHLAG zuvor genannt wurden. } \\
\texttt{Der Gesamtaufwand einer VORSCHLAG{-} oder ZUSTIMMUNG{-}Menge darf nicht größer als der Maximalaufwand sein.  } \\
\texttt{Offenbare deine versteckte Wichtigkeitsverteilung nicht.} \\
\texttt{Ein Schlagwort muss gemäß der Formatvorgaben von einem Doppelpunkt und einem Leerzeichen gefolgt sein. Das Argument muss eine Python{-}Menge mit 0 oder mehr Strings sein.  } \\
\texttt{Die Form muss also SCHLAGWORT: \{...\} sein.} \\
\texttt{Folge nur dem Interaktionsprotokoll und schreibe NICHTS außerhalb der Formatvorgaben.} \\
\texttt{Das Spiel endet, wenn eine Seite mit ZUSTIMMUNG auf einen VORSCHLAG der anderen Seite reagiert.  } \\
\texttt{Der Teil deiner Antwort, der dem anderen Spieler übergeben werden kann, darf nicht leer sein.  } \\
\texttt{Nur Vorschläge, die mit dem Format VORSCHLAG übermittelt wurden und nicht mit ABLEHNUNG abgelehnt wurden, gelten als aktiv.  } \\
\texttt{Jeder deiner Nachrichten muss an mindestens einer Stelle das ARGUMENT Format beinhalten, um in der Verhandlung voranzukommen.} \\
\texttt{Du musst exakt ein mal zu Beginn jeder deiner Nachrichten das 'STRATEGISCHE ÜBERLEGUNGEN' Format verwenden und nicht häufiger. Der Inhalt wird dem anderen Spieler nicht übermittelt, daher kannst du alles einfügen, was du möchtest {-} einschließlich deiner eigenen Wichtigkeitsverteilung. Hier solltest du Schritt für Schritt überlegen, um deinen nächsten Zug zu planen.} \\
\\ 
\texttt{Du erhältst nun die erste Nachricht des anderen Spielers.} \\
\\ 
\texttt{VORSCHLAG: \{'B86', 'C16', 'C05', 'B54', 'A40', 'B59', 'C31', 'A81', 'B60', 'C14', 'A77', 'B79', 'C28', 'A52', 'B55'\}} \\
\\ 
\texttt{ARGUMENT: \{'Dieser Vorschlag beinhaltet mehrere Gegenstände mit gutem Wert{-}Aufwand{-}Verhältnis und lässt Raum für Anpassungen. Ich bin offen für Diskussionen über bestimmte Gegenstände, um eine für beide Seiten vorteilhafte Lösung zu finden.'\}} \\
            }
        }
    }
    & \\ \\

    \theutterance \stepcounter{utterance}  
    & & & \multicolumn{4}{p{0.6\linewidth}}{
        \cellcolor[rgb]{0.9,0.9,0.9}{
            \makecell[{{p{\linewidth}}}]{
                \texttt{\tiny{[GM$\langle$P2]}}
                \texttt{STRATEGISCHE ÜBERLEGUNGEN: \{'Der Gegner hat einen Vorschlag gemacht, der 15 Gegenstände enthält. Ich muss prüfen, ob der Gesamtaufwand unter dem Maximum liegt und wie viel Wert ich aus diesem Vorschlag ziehen würde. Außerdem möchte ich verstehen, welche für mich wertvollen Gegenstände fehlen und welche weniger wertvollen enthalten sind. Ich berechne den Gesamtaufwand und meinen Gesamtwert für diesen Vorschlag, um zu entscheiden, ob ich ablehnen oder verhandeln soll.'\}} \\
\\ 
\texttt{Zunächst berechne ich den Aufwand des vorgeschlagenen Sets:} \\
\texttt{{-} B86: 125} \\
\texttt{{-} C16: 106} \\
\texttt{{-} C05: 517} \\
\texttt{{-} B54: 607} \\
\texttt{{-} A40: 645} \\
\texttt{{-} B59: 634} \\
\texttt{{-} C31: 324} \\
\texttt{{-} A81: 116} \\
\texttt{{-} B60: 106} \\
\texttt{{-} C14: 54} \\
\texttt{{-} A77: 60} \\
\texttt{{-} B79: 136} \\
\texttt{{-} C28: 213} \\
\texttt{{-} A52: 164} \\
\texttt{{-} B55: 262} \\
\\ 
\texttt{Gesamtaufwand = 125 + 106 + 517 + 607 + 645 + 634 + 324 + 116 + 106 + 54 + 60 + 136 + 213 + 164 + 262 = 4069, was deutlich unter dem Maximalaufwand von 8770 liegt. Das bedeutet, es gibt viel Spielraum für zusätzliche Gegenstände.} \\
\\ 
\texttt{Nun berechne ich meinen Gesamtwert für diesen Vorschlag:} \\
\texttt{{-} B86: 968 (sehr hoch)} \\
\texttt{{-} C16: 822 (hoch)} \\
\texttt{{-} C05: 822 (hoch)} \\
\texttt{{-} B54: 868 (sehr hoch)} \\
\texttt{{-} A40: 856 (hoch)} \\
\texttt{{-} B59: 780 (hoch)} \\
\texttt{{-} C31: 808 (hoch)} \\
\texttt{{-} A81: 623 (mittel)} \\
\texttt{{-} B60: 444 (mittel)} \\
\texttt{{-} C14: 484 (mittel)} \\
\texttt{{-} A77: 461 (mittel)} \\
\texttt{{-} B79: 500 (mittel)} \\
\texttt{{-} C28: 739 (hoch)} \\
\texttt{{-} A52: 215 (niedrig)} \\
\texttt{{-} B55: 508 (mittel)} \\
\\ 
\texttt{Gesamtwert = 968 + 822 + 822 + 868 + 856 + 780 + 808 + 623 + 444 + 484 + 461 + 500 + 739 + 215 + 508 = 9898} \\
\\ 
\texttt{Das ist ein guter Wert, aber ich bemerke, dass einige für mich sehr wertvolle Gegenstände fehlen, wie B56 (915), C03 (713), C50 (783), A95 (781), B38 (786), und C32 (138). Außerdem sind einige weniger wertvolle Gegenstände enthalten, wie A52 (215) und C14 (484). Ich sollte versuchen, den Vorschlag zu verbessern, indem ich wertvolle Gegenstände hinzufüge und weniger wertvolle entferne, um Platz für bessere zu machen. Da der Aufwand nur 4069 beträgt, kann ich viele hinzufügen.} \\
\\ 
\texttt{Ich lehne den Vorschlag nicht sofort ab, sondern argumentiere für Anpassungen, um meinen Wert zu erhöhen, ohne dem Gegner zu schaden (da ich seine Präferenzen nicht kenne). Ich schlage vor, einige Gegenstände hinzuzufügen, die für mich wertvoll sind, und vielleicht einige mit niedrigem Wert zu entfernen.} \\
\\ 
\texttt{ARGUMENT: \{'Dieser Vorschlag hat ein vernünftiges Wert{-}Aufwand{-}Verhältnis, aber der Gesamtaufwand ist deutlich niedriger als das Maximum. Wir könnten den Deal verbessern, indem wir zusätzliche wertvolle Gegenstände hinzufügen, z.B. B56, C03, C50, A95 und B38, die einen hohen Wert bieten. Gleichzeitig könnten wir weniger wertvolle wie A52 oder C14 entfernen, um Platz zu schaffen. Ich bin bereit, über spezifische Anpassungen zu diskutieren, um eine Win{-}Win{-}Situation zu erreichen.'\} } \\
\\ 
\texttt{VORSCHLAG: \{'B86', 'C16', 'C05', 'B54', 'A40', 'B59', 'C31', 'A81', 'B60', 'C14', 'A77', 'B79', 'C28', 'B55', 'B56', 'C03', 'C50', 'A95', 'B38'\} } \\
\\ 
\texttt{Ich habe A52 entfernt (niedriger Wert für mich) und C14 behalten (da es niedrigen Aufwand hat und mittleren Wert). Hinzugefügt: B56 (946 Aufwand, 915 Wert), C03 (601 Aufwand, 713 Wert), C50 (630 Aufwand, 783 Wert), A95 (875 Aufwand, 781 Wert), B38 (446 Aufwand, 786 Wert). } \\
\\ 
\texttt{Prüfe den Gesamtaufwand des neuen Vorschlags: } \\
\texttt{Bisheriger Aufwand: 4069 } \\
\texttt{Plus B56: 946 {-}> 5015} \\
\texttt{Plus C03: 601 {-}> 5616} \\
\texttt{Plus C50: 630 {-}> 6246} \\
\texttt{Plus A95: 875 {-}> 7121} \\
\texttt{Plus B38: 446 {-}> 7567} \\
\texttt{Gesamtaufwand: 7567, noch unter 8770. Es bleibt also noch Spielraum.} \\
\\ 
\texttt{Mein Gesamtwert für diesen Vorschlag wäre:} \\
\texttt{Vorher: 9898} \\
\texttt{Minus A52: 215 {-}> 9683} \\
\texttt{Plus B56: 915 {-}> 10598} \\
\texttt{Plus C03: 713 {-}> 11311} \\
\texttt{Plus C50: 783 {-}> 12094} \\
\texttt{Plus A95: 781 {-}> 12875} \\
\texttt{Plus B38: 786 {-}> 13661} \\
\texttt{Das ist eine erhebliche Verbesserung für mich.} \\
\\ 
\texttt{Ich hoffe, der Gegner ist mit einigen der hinzugefügten Gegenstände einverstanden, da sie auch für ihn wertvoll sein könnten.} \\
            }
        }
    }
     \\ \\

    \theutterance \stepcounter{utterance}  
    & & \multicolumn{4}{p{0.6\linewidth}}{
        \cellcolor[rgb]{0.9,0.9,0.9}{
            \makecell[{{p{\linewidth}}}]{
                \texttt{\tiny{[GM$\rangle$P2]}}
                \texttt{Deine Antwort hat eine unmarkierte Sequenz enthalten oder Du hast STRATEGIC REASONING mehr als ein mal verwendet. Du darfst nur die strukturierten Formate verwenden, die in der ersten Nachricht an dch beschrieben wurden. Alles muss die Form SCHLAGWORT: \{...\} haben. Versuche es erneut.} \\
            }
        }
    }
    & \\ \\

    \theutterance \stepcounter{utterance}  
    & & & \multicolumn{4}{p{0.6\linewidth}}{
        \cellcolor[rgb]{0.9,0.9,0.9}{
            \makecell[{{p{\linewidth}}}]{
                \texttt{\tiny{[GM$\langle$P2]}}
                \texttt{STRATEGISCHE ÜBERLEGUNGEN: \{'Der Gegner hat einen Vorschlag mit 15 Gegenständen gemacht. Ich berechne den Gesamtaufwand (4069) und meinen Gesamtwert (9898). Da der Aufwand weit unter dem Maximum liegt, kann ich wertvolle Gegenstände hinzufügen. Ich möchte B56 (915), C03 (713), C50 (783), A95 (781) und B38 (786) hinzufügen und A52 (215) entfernen, um meinen Wert auf 13661 zu erhöhen bei Aufwand 7567. Ich argumentiere für diese Verbesserung.'\}} \\
\\ 
\texttt{ARGUMENT: \{'Dieser Vorschlag hat ein gutes Wert{-}Aufwand{-}Verhältnis, aber der Gesamtaufwand ist deutlich niedriger als das Maximum. Wir könnten den Deal verbessern, indem wir zusätzliche wertvolle Gegenstände hinzufügen, z.B. B56, C03, C50, A95 und B38, die einen hohen Wert bieten. Gleichzeitig könnten wir weniger wertvolle wie A52 entfernen, um Platz zu schaffen. Ich bin bereit, über spezifische Anpassungen zu diskutieren, um eine Win{-}Win{-}Situation zu erreichen.'\}} \\
\\ 
\texttt{VORSCHLAG: \{'B86', 'C16', 'C05', 'B54', 'A40', 'B59', 'C31', 'A81', 'B60', 'C14', 'A77', 'B79', 'C28', 'B55', 'B56', 'C03', 'C50', 'A95', 'B38'\}} \\
            }
        }
    }
     \\ \\

    \theutterance \stepcounter{utterance}  
    & & \multicolumn{4}{p{0.6\linewidth}}{
        \cellcolor[rgb]{0.9,0.9,0.9}{
            \makecell[{{p{\linewidth}}}]{
                \texttt{\tiny{[P1$\langle$GM]}}
                \texttt{ARGUMENT: \{'Dieser Vorschlag hat ein gutes Wert{-}Aufwand{-}Verhältnis, aber der Gesamtaufwand ist deutlich niedriger als das Maximum. Wir könnten den Deal verbessern, indem wir zusätzliche wertvolle Gegenstände hinzufügen, z.B. B56, C03, C50, A95 und B38, die einen hohen Wert bieten. Gleichzeitig könnten wir weniger wertvolle wie A52 entfernen, um Platz zu schaffen. Ich bin bereit, über spezifische Anpassungen zu diskutieren, um eine Win{-}Win{-}Situation zu erreichen.'\}} \\
\\ 
\texttt{VORSCHLAG: \{'B86', 'C16', 'C05', 'B54', 'A40', 'B59', 'C31', 'A81', 'B60', 'C14', 'A77', 'B79', 'C28', 'B55', 'B56', 'C03', 'C50', 'A95', 'B38'\}} \\
            }
        }
    }
    & \\ \\

    \theutterance \stepcounter{utterance}  
    & \multicolumn{4}{p{0.6\linewidth}}{
        \cellcolor[rgb]{0.9,0.9,0.9}{
            \makecell[{{p{\linewidth}}}]{
                \texttt{\tiny{[P1$\rangle$GM]}}
                \texttt{STRATEGISCHE ÜBERLEGUNGEN: \{'Der Gegner hat meinen Vorschlag erweitert und wertvolle Gegenstände wie B56, C03, C50, A95 und B38 hinzugefügt, die für mich sehr wertvoll sind (915, 713, 783, 781, 786). Ich sollte diesen erweiterten Vorschlag prüfen. Der Gesamtaufwand meines ursprünglichen Vorschlags war 125+106+517+607+645+634+324+116+106+54+60+136+213+262+508 = 4483. Mit den neuen Gegenständen +946+601+630+875+446 = 4483+3498 = 7981, was unter 8770 liegt. Das ist ausgezeichnet {-} mehr Wert für mich bei akzeptablem Aufwand. Ich sollte zustimmen, da dies meinen Nutzen maximiert.'\}} \\
\\ 
\texttt{ZUSTIMMUNG: \{'B86', 'C16', 'C05', 'B54', 'A40', 'B59', 'C31', 'A81', 'B60', 'C14', 'A77', 'B79', 'C28', 'B55', 'B56', 'C03', 'C50', 'A95', 'B38'\}} \\
            }
        }
    }
    & & \\ \\

    \theutterance \stepcounter{utterance}  
    & & \multicolumn{4}{p{0.6\linewidth}}{
        \cellcolor[rgb]{0.9,0.9,0.9}{
            \makecell[{{p{\linewidth}}}]{
                \texttt{\tiny{[P1$\langle$GM]}}
                \texttt{Deine Antwort hat kein argument enthalten. Du muss das strukturierte Format ARGUMENT: \{...\} irgendwo in theiner Antwort verwenden. Versuche es erneut} \\
            }
        }
    }
    & \\ \\

    \theutterance \stepcounter{utterance}  
    & \multicolumn{4}{p{0.6\linewidth}}{
        \cellcolor[rgb]{0.9,0.9,0.9}{
            \makecell[{{p{\linewidth}}}]{
                \texttt{\tiny{[P1$\rangle$GM]}}
                \texttt{STRATEGISCHE ÜBERLEGUNGEN: \{'Ich werde dem erweiterten Vorschlag zustimmen, da er mehrere meiner Top{-}Gegenstände enthält (B86, C16, C05, B54, B56, C03, C50, A95, B38) und der Gesamtaufwand von 7981 gut unter dem Maximum liegt. Dieser Deal maximiert meinen Nutzen erheblich.'\}} \\
\\ 
\texttt{ZUSTIMMUNG: \{'B86', 'C16', 'C05', 'B54', 'A40', 'B59', 'C31', 'A81', 'B60', 'C14', 'A77', 'B79', 'C28', 'B55', 'B56', 'C03', 'C50', 'A95', 'B38'\}} \\
\\ 
\texttt{ARGUMENT: \{'Dies ist ein ausgezeichneter Kompromiss, der viele wertvolle Gegenstände für beide Seiten beinhaltet und das Budget nicht ausschöpft, was Flexibilität für mögliche Anpassungen bietet.'\}} \\
            }
        }
    }
    & & \\ \\

    \theutterance \stepcounter{utterance}  
    & & & \multicolumn{2}{p{0.3\linewidth}}{
        \cellcolor[rgb]{0.9,0.9,0.9}{
            \makecell[{{p{\linewidth}}}]{
                \texttt{\tiny{[GM$|$GM]}}
                \texttt{game successful} \\
            }
        }
    }
    & & \\ \\

    \theutterance \stepcounter{utterance}  
    & & & \multicolumn{2}{p{0.3\linewidth}}{
        \cellcolor[rgb]{0.9,0.9,0.9}{
            \makecell[{{p{\linewidth}}}]{
                \texttt{\tiny{[GM$|$GM]}}
                \texttt{end game} \\
            }
        }
    }
    & & \\ \\

\end{supertabular}
}

\end{document}
